% ---------------------------------------------------------------------
%  PHỤ LỤC
% ---------------------------------------------------------------------
\chapter{Phụ lục}
\label{ch:phuluc}
\markboth{Phụ lục}{}

\section{Hướng dẫn cài đặt và chạy hệ thống}
\label{sec:huongdan-run}

Mã nguồn của hệ thống nằm trong thư mục \code{E:\textbackslash
AI-EXAMHUB-DAU}. Cấu hình thật được đọc từ các tệp môi trường không nên đưa
lên kho công khai. Để chạy stack phát triển bằng Docker Compose, tạo bản
cấu hình cục bộ từ ví dụ và khởi động các dịch vụ:

\begin{lstlisting}[language=bash, caption={Các lệnh khởi chạy stack phát triển.}]
cd E:\AI-EXAMHUB-DAU
docker compose up --build
\end{lstlisting}

Theo tệp Compose hiện có, frontend được công bố tại
\code{http://localhost:8080}, backend tại \code{http://localhost:3000/api}
và PostgreSQL tại cổng 5432. Docker Compose cho phép mô tả và khởi chạy
nhiều container từ một cấu hình \cite{DockerComposeDocs}. Các giá trị cổng
và mật khẩu development trong dự án chỉ phục vụ cài đặt cục bộ; phải đổi
trước khi triển khai có dữ liệu thật.

Để chạy từng phần, vào \code{backend/api}, cài phụ thuộc theo
\code{package-lock.json}, chuẩn bị PostgreSQL và cấu hình các biến môi
trường từ \code{.env.example}; sau đó chạy \code{npm run dev}. Frontend
được chạy trong thư mục \code{frontend/web} bằng \code{npm install} và
\code{npm run dev}; đặt \code{VITE\_API\_URL} trỏ tới backend đang chạy.
Các lệnh kiểm tra là \code{npm test} tại backend và
\code{npm run build} tại frontend. Nếu database schema chưa đầy đủ, cần
áp dụng các migration theo đúng thứ tự trong thư mục dự án. Không chạy
test bằng tài khoản hoặc database production.

\section{Tổng quan cơ sở dữ liệu}
\label{sec:db-overview}

Schema được khởi tạo từ \code{database/schema.sql} và bổ sung theo các
migration. Định danh chính dùng UUID; các nội dung linh hoạt như lựa chọn,
câu trả lời, danh sách ID hoặc metadata được lưu dạng JSONB; thời điểm
nghiệp vụ dùng \code{TIMESTAMPTZ}. Các khoá ngoại quy định vòng đời phụ
thuộc của bài nộp, kết quả, lượt thi và thành viên. Các quan hệ lõi được
minh hoạ ở Hình~\ref{fig:erd}; Bảng~\ref{tab:dulieu} trong
Chương~\ref{ch:thietke} là bản đồ thực thể ở mức khái quát.

\subsection{Quy ước đọc từ điển dữ liệu}
Trong các bảng dưới đây, PK là khoá chính; FK là khoá ngoại; UQ là giá trị
duy nhất; NN là bắt buộc; CK là ràng buộc miền/điều kiện. Kiểu JSONB hỗ trợ
dữ liệu JSON, nhưng không tự thay thế kiểm tra tính hợp lệ của cấu trúc
payload trong API. Một số cột được bổ sung qua migration hoặc
\code{ensureExamSchema()} lúc ứng dụng khởi động; vì vậy schema cuối cùng
đang chạy phải được đối chiếu với toàn bộ migration, không chỉ một file SQL.

\section{Từ điển dữ liệu}
\label{sec:data-dictionary}

\subsection{Bảng users -- tài khoản và vai trò}
Bảng \code{users} là gốc định danh của người dùng, giữ email đăng nhập,
thông tin hồ sơ và phạm vi quyền. Vai trò được giới hạn vào
\code{admin}, \code{teacher} hoặc \code{student}. Trường
\code{institutional\_id} hỗ trợ mã cán bộ/sinh viên và có chỉ mục duy nhất
riêng khi trường khác NULL. Không lưu mật khẩu rõ; API lưu giá trị băm ở
\code{password\_hash}.

\begin{table}[H]
\centering
\caption{Từ điển dữ liệu bảng users.}
\label{tab:dict-users}
\begin{tabularx}{\textwidth}{>{\ttfamily\color{daumaroon}}p{0.20\textwidth}p{0.17\textwidth}p{0.12\textwidth}X}
\toprule
\rowcolor{gray100}Trường & Kiểu & Ràng buộc & Ý nghĩa \\
\midrule
id & UUID & PK & Định danh tài khoản, mặc định sinh UUID. \\
email & VARCHAR(255) & NN, UQ & Địa chỉ đăng nhập, duy nhất. \\
password\_hash & TEXT & NN & Mật khẩu đã băm; không phải mật khẩu gốc. \\
full\_name & VARCHAR(255) & NN & Họ tên hiển thị. \\
role & VARCHAR(32) & NN, CK & Vai trò admin/teacher/student. \\
avatar\_url & TEXT & NULL & Đường dẫn hoặc nội dung ảnh đại diện theo API. \\
phone & VARCHAR(30) & NULL & Số điện thoại liên hệ. \\
institutional\_id & VARCHAR(64) & UQ, NULL & Mã định danh cơ sở đào tạo. \\
faculty & VARCHAR(255) & NULL & Khoa của người dùng. \\
department & VARCHAR(255) & NULL & Bộ môn hoặc đơn vị. \\
cohort & VARCHAR(64) & NULL & Khoá học. \\
class\_name & VARCHAR(128) & NULL & Lớp hành chính. \\
created\_at & TIMESTAMPTZ & NN & Thời điểm tạo tài khoản. \\
\bottomrule
\end{tabularx}
\end{table}

\subsection{Bảng questions -- ngân hàng câu hỏi}
Bảng \code{questions} giữ nội dung và metadata cần cho việc chọn câu vào đề.
Các trường liên quan loại câu, câu trả lời, mức nhận thức và nguồn có thể
được mở rộng qua migration. Lựa chọn và rubric có cấu trúc linh hoạt; nội
dung nguồn được lưu riêng ở \code{source\_evidence} để phục vụ rà soát và
truy nguyên. Câu hỏi có thể tồn tại độc lập với người tạo nếu tài khoản
gốc bị xoá.

\begin{table}[H]
\centering
\caption{Từ điển dữ liệu bảng questions.}
\label{tab:dict-questions}
\begin{tabularx}{\textwidth}{>{\ttfamily\color{daumaroon}}p{0.21\textwidth}p{0.17\textwidth}p{0.12\textwidth}X}
\toprule
\rowcolor{gray100}Trường & Kiểu & Ràng buộc & Ý nghĩa \\
\midrule
id & UUID & PK & Định danh câu hỏi. \\
subject & VARCHAR(255) & NN & Môn/chủ đề học thuật. \\
content & TEXT & NN & Lời dẫn hoặc nội dung câu hỏi. \\
choices & JSONB & NN & Danh sách lựa chọn; mặc định mảng rỗng. \\
answer\_index & INTEGER & NULL & Chỉ số lựa chọn đúng cho câu trắc nghiệm. \\
answer\_text & TEXT & NULL & Đáp án tham khảo dạng văn bản. \\
source\_evidence & TEXT & NULL & Trích đoạn nguồn làm căn cứ cho câu hỏi. \\
cognitive\_level & VARCHAR(32) & NULL & Mức nhận thức như understanding/application. \\
rubric & JSONB & NULL & Cấu trúc tiêu chí chấm tự luận nếu có. \\
difficulty & VARCHAR(16) & NN, CK & easy, medium hoặc hard; mặc định medium. \\
type & VARCHAR(32) & Schema mở rộng & Loại multipleChoice/essay; được bổ sung trong schema runtime/migration. \\
created\_by & UUID & FK, NULL & Tài khoản tạo; xoá tài khoản đặt NULL. \\
created\_at & TIMESTAMPTZ & NN & Thời điểm tạo câu hỏi. \\
\bottomrule
\end{tabularx}
\end{table}

\subsection{Bảng exams -- đề thi}
Bảng \code{exams} lưu thông tin định danh và trạng thái vòng đời của đề.
Các chi tiết cấu hình đề được lưu trong \code{metadata} dạng JSONB; các
liên kết câu hỏi có thứ tự được chuẩn hoá ở bảng \code{exam\_questions}.
Schema cũng có \code{question\_ids} nhằm tương thích với các luồng cũ.
Khi bảo trì cần xác định một nguồn dữ liệu chuẩn để tránh hai biểu diễn
bị lệch nhau.

\begin{table}[H]
\centering
\caption{Từ điển dữ liệu bảng exams.}
\label{tab:dict-exams}
\begin{tabularx}{\textwidth}{>{\ttfamily\color{daumaroon}}p{0.21\textwidth}p{0.17\textwidth}p{0.12\textwidth}X}
\toprule
\rowcolor{gray100}Trường & Kiểu & Ràng buộc & Ý nghĩa \\
\midrule
id & UUID & PK & Định danh đề. \\
title & VARCHAR(255) & NN & Tên đề thi. \\
subject & VARCHAR(255) & NN & Môn học. \\
status & VARCHAR(16) & NN, CK & Trạng thái phát hành ở schema. \\
workflow\_status & VARCHAR(16) & NN/default & Trạng thái quy trình: draft/review/approved/published trong ứng dụng. \\
question\_ids & JSONB & NN & Danh sách ID tương thích với luồng legacy. \\
metadata & JSONB & NN & Thời lượng, tổng điểm, ma trận, năm học, rubric và thông tin đề. \\
created\_by & UUID & FK, NULL & Tài khoản tạo; xoá người tạo đặt NULL. \\
created\_at & TIMESTAMPTZ & NN & Thời điểm tạo. \\
updated\_at & TIMESTAMPTZ & NN & Thời điểm cập nhật gần nhất. \\
\bottomrule
\end{tabularx}
\end{table}

\subsection{Bảng exam\_questions -- câu hỏi thuộc đề}
Bảng nối biểu diễn quan hệ nhiều--nhiều giữa đề và câu hỏi, đồng thời
lưu thứ tự hiển thị. Một câu hỏi có thể tái sử dụng trong nhiều đề; xoá đề
sẽ cascade xoá liên kết nhưng không xoá câu hỏi.
\begin{table}[H]
\centering
\caption{Từ điển dữ liệu bảng exam\_questions.}
\label{tab:dict-exam-questions}
\begin{tabularx}{\textwidth}{>{\ttfamily\color{daumaroon}}p{0.21\textwidth}p{0.17\textwidth}p{0.12\textwidth}X}
\toprule
\rowcolor{gray100}Trường & Kiểu & Ràng buộc & Ý nghĩa \\
\midrule
exam\_id & UUID & PK, FK & Đề chứa câu hỏi; xoá đề cascade liên kết. \\
question\_id & UUID & PK, FK & Câu hỏi; không cho xoá nếu vẫn thuộc đề. \\
position & INTEGER & NN & Vị trí câu trong đề. \\
\bottomrule
\end{tabularx}
\end{table}

\subsection{Bảng exam\_rooms -- phòng thi}
Phòng thi gắn với một đề chính và có thể gắn nhiều đề bổ sung qua bảng
\code{exam\_room\_exams}. Mã tham gia được đảm bảo duy nhất. Thời điểm mở,
đóng và trạng thái điều khiển khả năng vào phòng.
\begin{table}[H]
\centering
\caption{Từ điển dữ liệu bảng exam\_rooms.}
\label{tab:dict-exam-rooms}
\begin{tabularx}{\textwidth}{>{\ttfamily\color{daumaroon}}p{0.21\textwidth}p{0.17\textwidth}p{0.12\textwidth}X}
\toprule
\rowcolor{gray100}Trường & Kiểu & Ràng buộc & Ý nghĩa \\
\midrule
id & UUID & PK & Định danh phòng thi. \\
exam\_id & UUID & FK, NN & Đề chính của phòng; xoá đề cascade phòng. \\
name & VARCHAR(255) & NN & Tên phòng. \\
starts\_at & TIMESTAMPTZ & NULL & Thời điểm dự kiến bắt đầu. \\
ends\_at & TIMESTAMPTZ & NULL & Thời điểm kết thúc/đóng nhận bài. \\
status & VARCHAR(16) & NN, CK & scheduled, active hoặc closed. \\
access\_code & VARCHAR(8) & NN, UQ & Mã tham gia ngẫu nhiên duy nhất. \\
created\_by & UUID & FK, NULL & Tài khoản tạo phòng. \\
\bottomrule
\end{tabularx}
\end{table}

\subsection{Bảng exam\_room\_exams -- đề có thể phân cho phòng}
Bảng liên kết phòng với một hoặc nhiều đề được giảng viên chọn. Mỗi đề có
vị trí trong danh sách; một lượt thi chọn đề được gán cho sinh viên và
lưu ID đề ở \code{exam\_attempts}.
\begin{table}[H]
\centering
\caption{Từ điển dữ liệu bảng exam\_room\_exams.}
\label{tab:dict-room-exams}
\begin{tabularx}{\textwidth}{>{\ttfamily\color{daumaroon}}p{0.21\textwidth}p{0.17\textwidth}p{0.12\textwidth}X}
\toprule
\rowcolor{gray100}Trường & Kiểu & Ràng buộc & Ý nghĩa \\
\midrule
room\_id & UUID & PK, FK & Phòng thi; xoá phòng cascade liên kết. \\
exam\_id & UUID & PK, FK & Đề được đưa vào danh sách phân bổ. \\
position & INTEGER & NN, UQ theo phòng & Vị trí đề; duy nhất trong một phòng. \\
\bottomrule
\end{tabularx}
\end{table}

\subsection{Bảng exam\_attempts -- lượt dự thi}
Lượt thi lưu trạng thái làm bài độc lập với định nghĩa đề. Bộ câu hỏi được
giao và đáp án lưu ở JSONB để tiếp tục bài làm; khoá duy nhất
\code{(room\_id, student\_id)} bảo đảm một lượt chính cho sinh viên trong
phòng theo schema hiện tại.
\begin{table}[H]
\centering
\caption{Từ điển dữ liệu bảng exam\_attempts.}
\label{tab:dict-attempts}
\begin{tabularx}{\textwidth}{>{\ttfamily\color{daumaroon}}p{0.21\textwidth}p{0.17\textwidth}p{0.12\textwidth}X}
\toprule
\rowcolor{gray100}Trường & Kiểu & Ràng buộc & Ý nghĩa \\
\midrule
id & UUID & PK & Định danh lượt thi. \\
room\_id & UUID & FK, NN & Phòng thi chứa lượt. \\
exam\_id & UUID & FK, NN & Đề được giao. \\
student\_id & UUID & FK, NN & Sinh viên; xoá tài khoản cascade lượt. \\
question\_set & JSONB & NN & Snapshot câu hỏi an toàn của lượt thi. \\
answers & JSONB & NN & Câu trả lời theo ID; mặc định object rỗng. \\
marked\_questions & JSONB & NN & Danh sách câu được đánh dấu xem lại. \\
current\_question & INTEGER & NN & Vị trí hiện tại để tiếp tục bài. \\
started\_at & TIMESTAMPTZ & NN & Thời điểm bắt đầu. \\
deadline\_at & TIMESTAMPTZ & NN & Hạn hoàn thành bài. \\
submitted\_at & TIMESTAMPTZ & NULL & Thời điểm nộp bài. \\
submission\_id & UUID & FK, UQ, NULL & Bài nộp đã tạo cho lượt thi. \\
updated\_at & TIMESTAMPTZ & NN & Thời điểm cập nhật gần nhất. \\
\bottomrule
\end{tabularx}
\end{table}

\subsection{Bảng exam\_attempt\_events -- sự kiện lượt thi}
Bảng sự kiện hỗ trợ giám sát một số dấu hiệu thao tác trình duyệt. Schema
hiện tại định nghĩa loại sự kiện \code{tab\_hidden}; đây không phải bằng
chứng đầy đủ về hành vi của sinh viên, cũng không thay thế cơ chế giám sát
trực tiếp hoặc quy chế thi.
\begin{table}[H]
\centering
\caption{Từ điển dữ liệu bảng exam\_attempt\_events.}
\label{tab:dict-attempt-events}
\begin{tabularx}{\textwidth}{>{\ttfamily\color{daumaroon}}p{0.21\textwidth}p{0.17\textwidth}p{0.12\textwidth}X}
\toprule
\rowcolor{gray100}Trường & Kiểu & Ràng buộc & Ý nghĩa \\
\midrule
id & UUID & PK & Định danh sự kiện. \\
attempt\_id & UUID & FK, NN & Lượt thi; xoá lượt cascade sự kiện. \\
event\_type & VARCHAR(32) & NN, CK & Loại sự kiện cho phép, hiện gồm tab\_hidden. \\
occurred\_at & TIMESTAMPTZ & NN & Thời điểm ghi nhận sự kiện. \\
\bottomrule
\end{tabularx}
\end{table}

\subsection{Bảng submissions -- bài nộp}
Submission là bản ghi nộp bài có cấu trúc thống nhất cho lịch sử và chấm
điểm. Câu trả lời được giữ trong JSONB; trạng thái phân biệt bài đã nộp,
đang chấm hoặc đã chấm. Sinh viên có thể bị đặt NULL khi tài khoản bị
xoá theo chính sách schema.
\begin{table}[H]
\centering
\caption{Từ điển dữ liệu bảng submissions.}
\label{tab:dict-submissions}
\begin{tabularx}{\textwidth}{>{\ttfamily\color{daumaroon}}p{0.21\textwidth}p{0.17\textwidth}p{0.12\textwidth}X}
\toprule
\rowcolor{gray100}Trường & Kiểu & Ràng buộc & Ý nghĩa \\
\midrule
id & UUID & PK & Định danh bài nộp. \\
exam\_id & UUID & FK, NN & Đề được nộp. \\
student\_id & UUID & FK, NULL & Sinh viên nộp. \\
answers & JSONB & NN & Nội dung trả lời đã nộp. \\
status & VARCHAR(16) & NN, CK & submitted, grading hoặc graded. \\
submitted\_at & TIMESTAMPTZ & NN & Thời điểm nộp. \\
\bottomrule
\end{tabularx}
\end{table}

\subsection{Bảng results -- kết quả chấm}
Mỗi bài nộp có tối đa một bản ghi kết quả. Kết quả giữ điểm, phản hồi và
thời điểm chấm. Mức điểm tối đa được xác định từ metadata của đề và được
kiểm tra trong luồng chấm.
\begin{table}[H]
\centering
\caption{Từ điển dữ liệu bảng results.}
\label{tab:dict-results}
\begin{tabularx}{\textwidth}{>{\ttfamily\color{daumaroon}}p{0.21\textwidth}p{0.17\textwidth}p{0.12\textwidth}X}
\toprule
\rowcolor{gray100}Trường & Kiểu & Ràng buộc & Ý nghĩa \\
\midrule
id & UUID & PK & Định danh kết quả. \\
submission\_id & UUID & FK, NN, UQ & Bài nộp; một kết quả cho mỗi bài. \\
score & NUMERIC(5,2) & NN & Điểm số, mặc định 0. \\
feedback & TEXT & NULL & Nhận xét của người chấm. \\
graded\_at & TIMESTAMPTZ & NN & Thời điểm ghi/cập nhật điểm. \\
\bottomrule
\end{tabularx}
\end{table}

\subsection{Bảng rubrics -- tiêu chí chấm}
Bảng \code{rubrics} gắn bộ tiêu chí với đề thi. Tiêu chí có dạng JSONB để
mô tả cấu trúc rubric mà không tạo bảng con cho từng cấp điểm. Việc lưu
rubric không đồng nghĩa đã có công cụ AI tự động áp dụng rubric.
\begin{table}[H]
\centering
\caption{Từ điển dữ liệu bảng rubrics.}
\label{tab:dict-rubrics}
\begin{tabularx}{\textwidth}{>{\ttfamily\color{daumaroon}}p{0.21\textwidth}p{0.17\textwidth}p{0.12\textwidth}X}
\toprule
\rowcolor{gray100}Trường & Kiểu & Ràng buộc & Ý nghĩa \\
\midrule
id & UUID & PK & Định danh bộ rubric. \\
exam\_id & UUID & FK, NN & Đề áp dụng; xoá đề cascade rubric. \\
title & VARCHAR(255) & NN & Tên tiêu chí chấm. \\
criteria & JSONB & NN & Danh sách tiêu chí/mức điểm; mặc định mảng rỗng. \\
\bottomrule
\end{tabularx}
\end{table}

\subsection{Bảng classroom\_rooms -- lớp học trực tuyến}
Lớp học lưu thông tin buổi học, chế độ kết nối và quyền nhận thành viên.
Mã tham gia là duy nhất; tài khoản tạo lớp là bắt buộc. Google Meet URL
được kiểm tra ở API khi chế độ kết nối là Meet.
\begin{table}[H]
\centering
\caption{Từ điển dữ liệu bảng classroom\_rooms.}
\label{tab:dict-classroom-rooms}
\begin{tabularx}{\textwidth}{>{\ttfamily\color{daumaroon}}p{0.21\textwidth}p{0.17\textwidth}p{0.12\textwidth}X}
\toprule
\rowcolor{gray100}Trường & Kiểu & Ràng buộc & Ý nghĩa \\
\midrule
id & UUID & PK & Định danh lớp. \\
name & VARCHAR(160) & NN & Tên lớp học. \\
description & TEXT & NN & Mô tả, mặc định chuỗi rỗng. \\
mode & VARCHAR(20) & NN, CK & google\_meet hoặc direct. \\
meet\_url & TEXT & NULL & Liên kết Google Meet nếu chế độ tương ứng. \\
access\_code & VARCHAR(8) & NN, UQ & Mã tham gia lớp duy nhất. \\
status & VARCHAR(16) & NN, CK & scheduled, live hoặc ended. \\
admissions\_open & BOOLEAN & NN & Có nhận sinh viên mới hay không. \\
started\_at & TIMESTAMPTZ & NULL & Thời điểm bắt đầu buổi học. \\
ended\_at & TIMESTAMPTZ & NULL & Thời điểm kết thúc buổi học. \\
created\_by & UUID & FK, NN & Giảng viên/quản trị viên tạo lớp. \\
created\_at & TIMESTAMPTZ & NN & Thời điểm tạo. \\
\bottomrule
\end{tabularx}
\end{table}

\subsection{Bảng classroom\_members -- thành viên lớp}
Bảng nối thành viên có khoá ghép để cùng một sinh viên không gia nhập
trùng một lớp. Lịch sử gia nhập được lưu theo thời gian.
\begin{table}[H]
\centering
\caption{Từ điển dữ liệu bảng classroom\_members.}
\label{tab:dict-classroom-members}
\begin{tabularx}{\textwidth}{>{\ttfamily\color{daumaroon}}p{0.21\textwidth}p{0.17\textwidth}p{0.12\textwidth}X}
\toprule
\rowcolor{gray100}Trường & Kiểu & Ràng buộc & Ý nghĩa \\
\midrule
classroom\_id & UUID & PK, FK & Lớp học; xoá lớp cascade thành viên. \\
student\_id & UUID & PK, FK & Tài khoản sinh viên; xoá tài khoản cascade quan hệ. \\
joined\_at & TIMESTAMPTZ & NN & Thời điểm tham gia. \\
\bottomrule
\end{tabularx}
\end{table}

\section{Danh mục API theo nhóm nghiệp vụ}
\label{sec:api-catalog}

Các endpoint dưới đây được tổng hợp từ route và tài liệu API trong dự án.
API có tiền tố \code{/api}; các endpoint đánh dấu “cần xác thực” yêu cầu
JWT, còn thao tác quản trị/giảng viên có kiểm tra vai trò. Hợp đồng HTTP
có thể được mô tả bằng OpenAPI để tạo tài liệu và kiểm thử nhất quán
\cite{OpenAPISpec}.

\subsection{Xác thực, tài khoản và ngân hàng câu hỏi}
\begin{table}[H]
\centering
\caption{Endpoint xác thực, người dùng và câu hỏi.}
\label{tab:api-auth-questions}
\begin{tabularx}{\textwidth}{>{\ttfamily\color{daumaroon}}p{0.13\textwidth}p{0.39\textwidth}X}
\toprule
\rowcolor{gray100}Method & Endpoint & Mục đích / quyền \\
\midrule
POST & \code{/auth/register} & Đăng ký tài khoản sinh viên. \\
POST & \code{/auth/login} & Đăng nhập bằng email/mật khẩu. \\
GET & \code{/auth/me} & Lấy hồ sơ/token mới; cần xác thực. \\
PUT & \code{/auth/profile} & Cập nhật hồ sơ tài khoản; cần xác thực. \\
PUT & \code{/auth/password} & Đổi mật khẩu; cần xác thực. \\
GET & \code{/auth/google} & Khởi tạo đăng nhập Google nếu được cấu hình. \\
GET & \code{/auth/google/callback} & Nhận kết quả chuyển hướng OAuth. \\
GET & \code{/users} & Danh sách tài khoản/phân trang; admin. \\
POST & \code{/users} & Tạo tài khoản quản trị; admin. \\
GET & \code{/questions} & Tra cứu, lọc và phân trang câu hỏi. \\
POST/PUT/DELETE & \code{/questions} & Thêm/sửa/xoá câu hỏi; teacher/admin, xoá/sửa còn kiểm tra quyền sở hữu. \\
\bottomrule
\end{tabularx}
\end{table}

\subsection{Tài liệu và quản lý đề thi}
\begin{table}[H]
\centering
\caption{Endpoint xử lý tài liệu và vòng đời đề thi.}
\label{tab:api-exams}
\begin{tabularx}{\textwidth}{>{\ttfamily\color{daumaroon}}p{0.13\textwidth}p{0.42\textwidth}X}
\toprule
\rowcolor{gray100}Method & Endpoint & Mục đích / quyền \\
\midrule
POST & \code{/documents/upload} & Tải và trích xuất tài liệu; teacher/admin. \\
POST & \code{/documents/analyze} & Phân tích nội dung và trả ứng viên câu hỏi; teacher/admin. \\
POST & \code{/exams/generate} & Tạo đề từ câu hỏi đã chọn; teacher/admin. \\
GET & \code{/exams} & Danh sách đề theo phạm vi người dùng. \\
GET & \code{/exams/:id} & Chi tiết đề; câu trả về cho sinh viên ẩn đáp án nội bộ. \\
PUT & \code{/exams/:id} & Cập nhật metadata/trạng thái đề. \\
DELETE & \code{/exams/:id} & Xoá đề theo quyền. \\
POST & \code{/exams/:id/variants} & Tạo mã đề và ánh xạ lại đáp án. \\
POST & \code{/exams/:id/validate} & Kiểm tra cấu trúc/ma trận đề. \\
\bottomrule
\end{tabularx}
\end{table}

\subsection{Phòng thi, bài nộp và lớp học}
\begin{table}[H]
\centering
\caption{Endpoint tổ chức thi, chấm bài và quản lý lớp học.}
\label{tab:api-rooms-classrooms}
\begin{tabularx}{\textwidth}{>{\ttfamily\color{daumaroon}}p{0.13\textwidth}p{0.43\textwidth}X}
\toprule
\rowcolor{gray100}Method & Endpoint & Mục đích / quyền \\
\midrule
POST & \code{/exam-rooms} & Teacher/admin tạo phòng và mã tham gia. \\
GET & \code{/exam-rooms} & Teacher/admin liệt kê phòng được quản lý. \\
PATCH & \code{/exam-rooms/:id/status} & Mở/đóng phòng; kiểm tra chủ sở hữu. \\
GET & \code{/exam-rooms/:id/attempts} & Giám sát tiến độ và sự kiện; teacher/admin. \\
GET & \code{/exam-rooms/:id/grading} & Lấy các lượt thi đã nộp để chấm. \\
POST & \code{/exam-rooms/:code/start} & Sinh viên bắt đầu/tiếp tục bài bằng mã phòng. \\
PUT & \code{/exam-rooms/attempts/:id/answers} & Lưu câu trả lời và vị trí; chủ lượt thi. \\
POST & \code{/exam-rooms/attempts/:id/submit} & Nộp bài, tự chấm trắc nghiệm. \\
POST & \code{/exam-rooms/attempts/:id/grade} & Teacher/admin lưu điểm/phản hồi tự luận. \\
GET & \code{/exam-rooms/my-attempts} & Sinh viên xem lịch sử dự thi. \\
POST & \code{/classrooms} & Teacher/admin tạo lớp học. \\
POST & \code{/classrooms/join} & Sinh viên gia nhập bằng mã lớp. \\
GET/PATCH & \code{/classrooms/:id} & Xem lớp hoặc đổi trạng thái/nhận thành viên theo quyền. \\
\bottomrule
\end{tabularx}
\end{table}

\section{Danh mục migration dữ liệu}
\label{sec:migrations}

Các migration trong dự án ghi nhận việc mở rộng schema theo chức năng.
Bảng~\ref{tab:migrations} tóm tắt từ tên file; trước khi chạy trên dữ liệu
thật cần đọc nội dung migration, xác minh thứ tự, sao lưu và có phương án
rollback.
\begin{table}[H]
\centering
\caption{Danh mục migration hiện diện trong dự án.}
\label{tab:migrations}
\begin{tabularx}{\textwidth}{>{\ttfamily\color{daumaroon}}p{0.16\textwidth}p{0.47\textwidth}X}
\toprule
\rowcolor{gray100}Migration & Phạm vi thể hiện ở tên/tệp & Nhóm dữ liệu \\
\midrule
001 & \code{initial\_schema} & Schema ban đầu: người dùng, câu hỏi, đề và dữ liệu liên quan. \\
002 & \code{add\_metadata\_to\_exams} & Metadata cấu hình cho đề thi. \\
003 & \code{exam\_generation\_validation} & Trường phục vụ kiểm tra sinh đề. \\
004 & \code{add\_answer\_text\_to\_questions} & Đáp án dạng văn bản cho câu hỏi. \\
005 & \code{add\_question\_source\_evidence} & Bằng chứng trích từ tài liệu nguồn. \\
006 & \code{add\_essay\_cognitive\_rubric} & Mức độ nhận thức và rubric câu tự luận. \\
007 & \code{online\_exam\_rooms} & Phòng thi, lượt thi trực tuyến và sự kiện. \\
008 & \code{user\_profiles} & Trường hồ sơ tài khoản. \\
009 & \code{university\_user\_profiles} & Mã và thuộc tính hồ sơ đại học. \\
010 & \code{exam\_tab\_switch\_events} & Lịch sử sự kiện chuyển tab. \\
011 & \code{classrooms} & Lớp học, thành viên và mã tham gia. \\
012 & \code{classroom\_lifecycle} & Trạng thái vòng đời và nhận thành viên lớp. \\
\bottomrule
\end{tabularx}
\end{table}

\section{Ma trận kiểm thử bổ sung}
\label{sec:test-matrix}

Bên cạnh 35 ca tự động đã chạy, Bảng~\ref{tab:manual-tests} là bộ kịch bản
kiểm thử thủ công đề xuất cho buổi trình diễn. Đây không phải kết quả đã
chạy; cột trạng thái được để trống để người nộp cập nhật sau khi thực hiện.
\begin{table}[H]
\centering
\caption{Kịch bản kiểm thử thủ công trước khi bảo vệ.}
\label{tab:manual-tests}
\begin{tabularx}{\textwidth}{>{\bfseries\color{daumaroon}}p{0.10\textwidth}X X p{0.15\textwidth}}
\toprule
\rowcolor{gray100}Mã & Bước kiểm thử & Kết quả mong đợi & Trạng thái \\
\midrule
M01 & Đăng nhập lần lượt bằng vai trò sinh viên và giảng viên. & Menu và endpoint phù hợp vai trò; quyền không hợp lệ bị từ chối. & Chưa ghi nhận \\
M02 & Tải PDF có văn bản thật và kiểm tra phần text trích xuất. & Nội dung có thể xem trước; câu hỏi liên kết được với đoạn nguồn. & Chưa ghi nhận \\
M03 & Tải PDF quét chất lượng thấp hoặc file rỗng. & Có cảnh báo trích xuất; không sinh đề thành công giả. & Chưa ghi nhận \\
M04 & Tạo đề ma trận vượt số câu hiện có. & Hiện rõ số câu thiếu; không phát hành đề không hợp lệ. & Chưa ghi nhận \\
M05 & Gửi duyệt và chuyển trạng thái đề. & Trạng thái lưu sau khi tải lại trang. & Chưa ghi nhận \\
M06 & Vào phòng thi, trả lời một phần, tải lại rồi tiếp tục. & Đáp án/vị trí được khôi phục từ đúng lượt. & Chưa ghi nhận \\
M07 & Nộp bài có trắc nghiệm và tự luận. & Trắc nghiệm được chấm; tự luận chờ giảng viên chấm. & Chưa ghi nhận \\
M08 & Sinh viên thử xem API chi tiết đề/đáp án nội bộ. & API không trả đáp án đúng của bài thi. & Chưa ghi nhận \\
M09 & Gia nhập lớp đã kết thúc hoặc đóng nhận thành viên. & Từ chối người mới; thành viên đã có vẫn theo chính sách. & Chưa ghi nhận \\
M10 & Mở giao diện trên thiết bị di động, bàn phím và phóng đại. & Bố cục và thao tác chính còn sử dụng được. & Chưa ghi nhận \\
\bottomrule
\end{tabularx}
\end{table}

\section{Rủi ro, giới hạn và biện pháp giảm thiểu}
\label{sec:risk-register}

\begin{table}[H]
\centering
\caption{Sổ đăng ký rủi ro của nguyên mẫu AI ExamHub DAU.}
\label{tab:risk-register}
\begin{tabularx}{\textwidth}{>{\bfseries\color{daumaroon}}p{0.17\textwidth}X X}
\toprule
\rowcolor{gray100}Rủi ro & Tác động & Biện pháp đề xuất \\
\midrule
Đáp án AI sai & Sinh viên nhận nội dung sai; đề thi mất độ tin cậy. & Bắt buộc giảng viên xác minh câu, đáp án và nguồn trước phát hành. \\
Prompt injection từ tài liệu & Nội dung tải lên có thể chứa chỉ dẫn làm lệch yêu cầu cho mô hình. & Cô lập tài liệu khỏi instruction, loại hành động công cụ nhạy cảm, kiểm tra đầu ra và quyền duyệt của người. \\
Lộ dữ liệu qua API LLM & Học liệu/tài liệu có dữ liệu cá nhân đi ra nhà cung cấp. & Cảnh báo trước khi gửi, lọc dữ liệu nhạy cảm, cấu hình chính sách lưu giữ và nhà cung cấp được phép. \\
Lộ đáp án thi & Lộ đáp án làm mất công bằng kiểm tra. & Backend chỉ trả question view đã loại answer fields; test truy cập trực tiếp bằng vai trò sinh viên. \\
Chiếm đoạt tài khoản & Truy cập hồ sơ, đề, kết quả và dữ liệu lớp. & Secret mạnh, HTTPS, xác thực phù hợp, giới hạn quyền, nhật ký và quy trình đổi/thu hồi token. \\
Tệp độc hại/quá lớn & Cạn bộ nhớ, lỗi parser hoặc ảnh hưởng dịch vụ. & Giới hạn kích thước/định dạng, kiểm tra nội dung, cách ly xử lý và giới hạn tài nguyên. \\
OCR tiếng Việt kém & Nội dung nguồn bị nhận sai rồi lan truyền sang câu hỏi. & Định cấu hình ngôn ngữ phù hợp, hiển thị nguồn, đo lỗi OCR và cho phép người dùng hiệu chỉnh. \\
Database development & Mất dữ liệu hoặc lộ mật khẩu do cấu hình mặc định. & Thay mật khẩu, tách network, sao lưu/khôi phục kiểm chứng và dùng secret manager. \\
\bottomrule
\end{tabularx}
\end{table}

Sổ đăng ký là phân tích rủi ro thiết kế, không phải báo cáo kiểm thử xâm
nhập. Các rủi ro prompt injection và phân quyền dịch vụ API cần được đối
chiếu với khuyến nghị OWASP về ứng dụng web, API và LLM
\cite{OWASPTop10,OWASPAPI,OWASPLLM}; khung NIST giúp tổ chức việc nhận
diện, đo và xử lý rủi ro AI theo vòng đời \cite{NISTAIRMF}.

\section{Mục tự kiểm trước khi nộp}
\label{sec:checklist}

\begin{tukiem}
  \item Đã thay tên, MSSV, lớp, giảng viên hướng dẫn và năm học trên bìa.
  \item Đã rà soát mọi chương, mục lục, danh mục hình/bảng và số tham chiếu.
  \item Đã kiểm tra lại nội dung PDF/DOCX/OCR, các hạn chế kỹ thuật và phiên
        bản code được dùng để báo cáo.
  \item Đã chạy lại API tests, frontend build và ghi đúng phiên bản kết quả.
  \item Đã kiểm tra các đề xuất AI với tài liệu gốc; không coi confidence là
        xác suất đúng.
  \item Đã bổ sung ảnh chụp giao diện thật nếu giảng viên yêu cầu; các hình
        wireframe trong báo cáo được nhận diện là hình minh hoạ.
  \item Đã kiểm tra và ghi nguồn cho tối thiểu 15 tài liệu tham khảo có liên
        kết; mọi tài liệu trích dẫn đều có trong danh mục.
  \item Đã thay toàn bộ cấu hình development trước khi triển khai có dữ liệu
        thật.
  \item Đã xác nhận số trang sau khi biên dịch XeLaTeX; không coi số trang
        ước lượng trong source là số trang PDF.
  \item Đã xoá nội dung hướng dẫn và placeholder không thuộc bản nộp cuối.
\end{tukiem}

\section{Nhận xét của giảng viên hướng dẫn}
\label{sec:nhanxet}

\khungnhanxet[7.5cm]

\vspace{0.8em}
\noindent\begin{minipage}[t]{0.5\textwidth}
\color{daunavy}
{\color{daumaroon}\bfseries Điểm số}

\vspace{0.7em}
\dongchamnhan{Bằng số:}

\vspace{0.9em}
\dongchamnhan{Bằng chữ:}
\end{minipage}\hfill
\okyten{\dauGVHD}{Giảng viên hướng dẫn}
